% Vorgaben Assignment aus Studienheft SQL03
% Formatvorgaben fuer den Text
% Umfang: 8 - 10 Seiten (inkl. Abbildungen und Tabellen, aber ohne Deckblatt, % Gliederung und Literaturverzeichnis, Eidesstattliche Erklaerung)
% Zeilenabstand: 1,5
% Schriftart: frei
% Schriftgrad: 12 pt
% Variablen, physikalische Groessen und Funktionszeichen werden kursiv gedruckt.
% Korrekturrand: links: 4,5 cm, rechts 2,0 cm, oben und unten jeweils 3,0 cm
% Deckblatt: (Adresse, AKAD-E-Mail-Adresse, Immatrikulationsnummer, Modul-
% bezeichnung, Thema, Datum, Felder für Korrektor)
% Gliederung (1 Seite)
% Literaturverzeichnis (3 - 5 Literaturquellen  z. B. Lehrbuecher, aktuelle Fachartikel recherchieren)
% Eidesstattliche Erklaerung (unterschrieben und fest eingebunden)
% Bearbeitungsdauer: 2 Monate


\documentclass[a4paper,12pt]{article}
\usepackage[ngerman]{babel}
\usepackage[nottoc]{tocbibind} % Anzeigen des Literaturverzeichnisses im TOC
\usepackage{epsfig}
\usepackage{times}
\usepackage{supertabular}
\usepackage{wrapfig}
\usepackage{multirow}
\usepackage[onehalfspacing]{setspace}
\usepackage{listings}
\usepackage{mathptmx}
\usepackage{geometry}
\usepackage{helvet}
\usepackage{courier}
\usepackage{setspace}
\usepackage{textcomp}
\usepackage[T1]{fontenc}
\usepackage[utf8]{inputenc}
\usepackage{fancyhdr}
\usepackage{float} % Notwendig fuer figure[h]
\usepackage[printonlyused]{acronym}

\newif\iflistoffigures
\newif\iflistoftables
\newif\ifacronym

%Titel
\newcommand*{\Titel}{Digitale Kontaktaufnahme im Endkundenbereich} 

%Betreff
\newcommand*{\Betreff}{Projektbericht} 

%Betreuer
\newcommand*{\Betreuer}{Prof. Dr. Ulrich Kreutle} 

%Vor- und Nachname
\newcommand*{\Name}{Stefan Waidele}

%Straße und Hausnummer
\newcommand*{\Strasse}{Ensisheimer Straße 2} 

%Plz und Ort
\newcommand*{\PlzOrt}{79395 Neuenburg am Rhein} 

%Immatrikulationsnummer
\newcommand*{\Immatrikulationsnummer}{102 81 71}

%Email 
\newcommand*{\Email}{Stefan@Waidele.info} 


% Verzeichnisse (Wenn nicht benötigt, Zeile mit % auskommentieren oder löschen

%% Abbildungsverzeichnis 
%\listoffigurestrue
%% Tabellenverzeichnis
%\listoftablestrue
%% Abkürzungsverzeichnis
%\acronymtrue


% Wittwen und Waisen verhindern
\clubpenalty10000
\widowpenalty10000
\displaywidowpenalty=10000

\usepackage[flushmargin,hang,ragged]{footmisc}
\usepackage{lmodern} %Type1-Schriftart für nicht-englische Texte
\usepackage{fancyhdr}

\usepackage[
	pdftitle={\Titel},
	pdfsubject={SWE02 -- Sicherheitsaspekte einer Web--Anwendung},
	pdfauthor={Stefan Waidele},
	pdfkeywords={akad, swe02, assignment, wirtschaftsinformatik}
	hyperfootnotes=false,
	colorlinks=true,
	linkcolor=black,
	urlcolor=black,
	citecolor=black
]{hyperref}

%\renewcommand{\familydefault}{\rmdefault}
\renewcommand{\bflabel}[1]{\normalfont{\normalsize{#1}}\hfill}

%% Definition for Codeschnipsel im Fließtext
\newcommand{\code}{\texttt}
\newcommand{\todo}[1]{\fbox{\parbox{\textwidth}{\textbf{To do:} #1}}}
%\newcommand{\myref}[1]{„\ref{#1}~\nameref{#1}“}
\newcommand{\myref}[1]{\textit{\ref{#1}~\nameref{#1}}}

%% Für Codeblöcke mit Syntax-Highlighting
%% http://www.ctan.org/tex-archive/macros/latex/contrib/minted/
\usepackage{minted}
\definecolor{bg}{rgb}{0.95,0.95,0.95}


\makeatother

\geometry{a4paper, left=45mm, right=20mm, top=30mm, bottom=30mm}
\pagenumbering{roman}

\begin{document}
	
\parskip=1em
\parindent=0cm

%description: Deckblatt in Deutsch
%% Basierend auf einer TeXnicCenter-Vorlage von Tino Weinkauf.
%% sowie der akad-vorlage von Daniel Falkner
%%%%%%%%%%%%%%%%%%%%%%%%%%%%%%%%%%%%%%%%%%%%%%%%%%%%%%%%%%%%%%

%%%%%%%%%%%%%%%%%%%%%%%%%%%%%%%%%%%%%%%%%%%%%%%%%%%%%%%%%%%%%
%% Deckblatt
%%%%%%%%%%%%%%%%%%%%%%%%%%%%%%%%%%%%%%%%%%%%%%%%%%%%%%%%%%%%%
%%
%% ACHTUNG: Sie benötigen ein Hauptdokument, um diese Datei
%%          benutzen zu können. Verwenden Sie im Hauptdokument
%%          den Befehl "\input{dateiname}", um diese
%%          Datei einzubinden.
%%

\begin{titlepage}

%\thispagestyle{empty}

\vspace{5cm}

\Name \\ 
\Strasse \\ 
\PlzOrt\\ 
\href{mailto:\Email}{\Email}

AKAD University\\
Immatrikulationsnummer: \Immatrikulationsnummer

\vfill

Projektbericht
\LARGE
\textsc{Möglichkeiten der digitalen Kontaktaufnahme im Endkundenbereich}\\

\vfill

\normalsize

Betreuer: 

\today %%Datum der Abgabe - am besten selbst reinschreiben.

\vfill

\includegraphics[width=3cm]{akad_logo.png}

AKAD University

\end{titlepage}


%\includegraphics[scale=0.35]{akad_logo.png}

%\clearpage
\normalsize


\normalsize

\begin{spacing}{1.0} % Verzeichnisse werden mit einzeiligem Abstand gesetzt
\parskip=0em
\newpage

% Inhaltsverzeichnis
\setcounter{tocdepth}{2}
\tableofcontents 
\newpage

% Abbildungsverzeichnis
\iflistoffigures
\listoffigures 
\newpage
\fi

% Tabellenverzeichnis
\iflistoftables
\listoftables 
\newpage
\fi

% Abkürzungsverzeichnis
\ifacronym
\section*{Abkürzungsverzeichnis}
\addcontentsline{toc}{section}{Abkürzungsverzeichnis} 
\begin{acronym}[\hspace{5cm}]
	\acro{CAPTCHA}{Completely Automated Public Tourint Test to Tell Computers and Humans Apart}
	\acro{CSS}{Cascading Style Sheets}
	\acro{DoS}{Denial of Service}
	\acro{DDoS}{Distributed Denial of Service}
	\acro{ERP}{Enterprise Resource Planning}
	\acro{HTML}{Hypertext Markup Language}
	\acro{HTTP}{Hypertext Tranfsfer Protocol}
	\acro{HTTPS}{Hypertext Tranfsfer Protocol Secure}
	\acro{OWASP}{Open Web Application Security Project}
	\acro{PDO}{PHP Data Objects}
	\acro{SCM}{Supply Chain Management}
	\acro{SQL}{Structured Query Language}
	\acro{WWW}{World Wide Web}
	\acro{XSS}{Cross Site Scripting}
%%
%	
%	\acro{}{}
%	\acro{}{}
%	\acro{}{}
%	\acro{}{}
%	\acro{}{}
%	\acro{}{}
%	\acro{}{}
%	\acro{}{}
	\acro{Test}{Wird nicht im Text verwendet und taucht auch nicht im Verzeichnis auf}
\end{acronym}

\fi

\parskip=1em
\end{spacing} 

\clearpage

\newcounter{romanPagenumber} 
\setcounter{romanPagenumber}{\value{page}} % Roemische Seitenanzahl speichern.


\pagestyle{fancy}
\fancyhead{}
\fancyhead[LO,RE]{\textsc{\Titel}}
\fancyhead[RO,LE]{\thepage}
\fancyfoot[CO,CE]{}
\setlength{\headheight}{15pt}

% \nocite{*} 

\pagenumbering{arabic}

\begin{spacing}{1.5} % Zeilenabstand: 1,5 fuer den Textteil

\section{Einleitung}
\label{sec:einleitung}

\subsection{Begründung der Problemstellung}

Im Februar 2014 nutzen 40,4 Millionen Deutsche ein Smartphone\footnote{vgl. \cite{netzoekonom}}. Dies entspricht einem Anteil von knapp 50\% der Bevölkerung\footnote{vgl. \cite{destatis:bev}}, bzw. durschscnittlich eines in fast jedem deutschen Haushalt\footnote{vgl. \cite{destatis:hh}}. 

Mit dieser stetig steigenden Verbreitung von Smartphones in der Bevölkerung nehmen auch die Möglichkeiten der digitalen Kontaktaufnahme zu. 
Hierbei konkurieren verschiedene Technologien miteinander. Allen gemeinsam ist es, dass die Kunden oder Interessenten vor Ort auf digitale, meist Onlineangebote hingewiesen werden.

Unternehmen, die diese Möglichkeiten nutzen möchten, brauchen eine verlässliche Basis für die Entscheidung welche Technologie zum Einsatz kommen sollen. Obwohl technisch nichts dagegen spricht, mehrere Technologien gleichzeitig anzubieten, so nehmen sich diese dann doch den Platz und die Aufmerksamkeit des Kunden.

Wegen des in kleinen Betrieben schon aufgrund der Umsatzhöhe beschränkten Budget sowohl für IT als auch für Marketing ist diese Entscheidung noch bedeutungsvoller als bei großen Konzernen. Konzeption und Umsetzung werden in diesen Betrieben meist von Mitarbeitern erledigt, die gleichzeitig noch andere Aufgaben zu erfüllen haben. Daher sind auch die Bewertungskriterien entsprechend diesen Erfordernissen zu gestalten.

Ein stark wachsender Anteil von 27\% der Reisenden informieren sich wärend ihrer Reise über das mobile Internet. Genze 36\% der Reisendenden  teilen Ihre Reiseerlebnisse im Social Web\footnote{vgl. \cite{reiseanalyse}, Seite 5}. Somit ist das Thema für touristische Betriebe ein Zukunftsthema, das auch schon heute relevant ist.

\subsection{Ziele dieser Arbeit}

Ziel dieses Projekts ist es, Handlungsempfehlungen für den Einsatz von Technologien zur digitalen Kontaktaufnahme in kleinen Betrieben der Tourismusbranche zu geben.

\subsection{Aufbau der Arbeit}

\todo{An die tatsächliche Arbeit anpassen}

Zunächst werden im Kapitel~\ref{sec:grundlagen} die für diese Arbeit relevanten Begriffe und Konzepte definiert, bevor im Kapitel~\ref{sec:technologien} die momentan verfügbaren Technologien genannt und erklärt werden. Die Erarbeitung von Kriterien für die Bewertung und Vergleich der Technologien in Kapitel~\ref{sec:kriterien} wird von zwei Befragungen unterstüzt, die in Anhang~\ref{sec:kundenbefragung} und Anhang~\ref{sec:betriebsbefragung} beschrieben sind. Hiermit schließt der theoretische Teil dieser Arbeit. 

Nach einer kurzen Vorstellung der Tourismusbranche in Kapitel~\ref{sec:tourismusbranche} werden die relevanten Technologien in Kapitel~\ref{sec:bewertung} bewertet und miteinander Verglichen.

In den Kapiteln~\ref{sec:einsatzzwecke} und \ref{sec:einsatzorte} werden dann Handlungsempfehlungen bezüglich möglichen Einsatzzwecken und Einsatzorten gegeben.

\subsection{Abgrenzung}

\todo{An die tatsächliche Arbeit anpassen}

Da der Schritt von der persönlichen  zur unpersönlichen Kommunikation getan wird, sollte untersucht werden, für welche Zielsetzungen dies in Betracht kommt und wünschenswert ist. Da dies den Umfang dieser Arbeit sprengen würde, wird diese Fragestellung nicht ausführlich behandelt.

„Einsatzmöglichkeiten III: Gestaltung“ wird nicht behandelt, hierfür wären spezielle Untersuchungen notwendig.

Usability wird nur anhand von groben Einteilungen behandelt und nicht detailliert untersucht.

Die Technologien Internet--Nachrichtendienste, Mobilfunk--Nachrichtendienste und Hotel––App werden nicht erschöpfend betrachtet werden, da sie aufgrund der Bewertungskriterien Kapital-- bzw. Personaleinsatz offensichtlich nicht optimal für kleine Betriebe sind.
\section{Grundlagen}
\label{sec:grundlagen}

\subsection{Digitale Kontaktaufnahme}

Definition als Vorgang vorhandene oder potentielle
Kunden vor Ort abzuholen und auf ein digitales bzw. Onlineangebot zu
führen.

\subsection{Eingrenzung der Kundengruppe}

Beschränkung der potentiellen und vorhandenen Kunden kleiner touristischer Betriebe.

\subsection{Notwendigkeit des Medienbruchs}

Erläuterung warum ein Wechsel des Kommunikationsmediums von direkt (bzw. Print) zum Digitalen (Online oder Offline) wünschenswert ist.

Customer Journey, Reiseanalyse 2014, S5, unten! und http://www.fur.de/ra/news-daten/aktueller-newsletter/nl-0714-die-ra-customer-journey/ 

\section{Technologische Möglichkeiten der digitalen Kontaktaufnahme}
\label{sec:technologien}

\subsection{QR–Codes}
\subsection{NFC--Tags}
\subsection{Bluetooth}
\subsection{Kurzlinks}
\subsection{Internet--Nachrichtendienste}

z.B. Twitter für nicht an eine Mobilfunknummer gebundene Nachrichtendienste. Auch: iMessage, evt. Facebook Messenger

Aufgrund der hohen Anforderung bzgl. syncroner Kommunikation ist diese Möglichkeit nicht optimal für kleine Betriebe und wird nur im Grundlagenteil behandelt.

\subsection{Mobilfunk--Nachrichtendienste}

z.B. WhatsApp für Nachrichtendienste, die an eine Mobilfunknummer gebunden sind, selbst wenn die eigentliche Kommunikation über Internettechniken abgewickelt wird. Auch: SMS.

Aufgrund der hohen Anforderung bzgl. syncroner Kommunikation ist diese Möglichkeit nicht optimal für kleine Betriebe und wird nur im Grundlagenteil behandelt.

\subsection{Hotel--Apps}

z.B. Protel „Voyager“

Aufgrund der hohen Anforderung bzgl. Kapitaleinsatz ist diese Möglichkeit nicht optimal für kleine Betriebe und wird nur im Grundlagenteil behandelt.

\section{Bewertungskriterien}
\label{sec:kriterien}

Hier wäre u.U. eine Befragung von Kunden und Betrieben bzgl. der Gewichtung der Kriterien (und zur Gewinnung von weiteren Kriterien) sinnvoll.

\subsection{Benötigte Hard-- und Software (Hotel)}
\subsection{Benötigte Hard-- und Software (Kunde)}
\subsection{Bekanntheitsgrad}
\subsection{Verbreitungsgrad}
\subsection{Benutzerfreundlichkeit}
\subsection{Technische Möglichkeiten}
z.B. Speicherkapazität

\subsection{Finanzielle Anforderungen}
\subsection{Personelle Anforderungen}
\subsection{Mißbrauchsmöglichkeiten}
z.B. überklebte QR--Codes
\section{Angriffsarten und Gegenmaßnahmen}
\label{sec:hauptteil}


\section{Fazit \& Ausblick}

\subsection{Fazit}


\subsection{Ausblick}


\appendix
\section{Befragung der Kunden}
\label{sec:kundenbefragung}
\subsection{Fragebogen}
\subsection{Ergebnisse}

\section{Befragung der Unternehmen}
\label{sec:betriebsbefragung}
\subsection{Fragebogen}
\subsection{Ergebnisse}


\end{spacing}

\clearpage

% Literaturverzeichniss - Ab hier wieder Roemische Seitenzahlen
\pagestyle{plain}
\pagenumbering{roman}
\setcounter{page}{\theromanPagenumber}
\bibliographystyle{apalike}
\bibliography{literatur}
\onehalfspacing
\clearpage

\pagestyle{empty} 
\thispagestyle{empty}

\begin{center}
{\Large Eidesstattliche Erklärung}
\vspace*{4cm}\end{center}
\noindent
Ich versichere, dass ich das beiliegende Assignment selbstständig verfasst, keine anderen als die angegebenen Quellen und Hilfsmittel benutzt sowie alle wörtlich oder sinngemäß übernommenen Stellen in der Arbeit gekennzeichnet habe. 
\vspace{3cm}

\hspace{-0.8cm}
\rule[0.5ex]{6.5cm}{1pt}
\hspace{1.3cm}
\rule[0.5ex]{6.5cm}{1pt}
\\(Datum, Ort)
\hspace{6.3cm}
(Unterschrift)

\clearpage

%Messbox zur Druckkontrolle:
\newcommand{\Messbox}[2]{% Parameters: #1=Breite, #2=Hoehe
\setlength{\unitlength}{1.0mm}%
\begin{picture}(#1,#2)%
\linethickness{0.05mm}%
\put(0,0){\dashbox{0.2}(#1,#2)%
{\parbox{#1mm}{%
\centering\footnotesize 
%{\bf MESSBOX}\\ 
Breite $ = #1 {\rm\ mm}$\\
H\"ohe $ = #2 {\rm\ mm}$
}}}\end{picture}
}

\begin{center}
\textbf{--- Druckgröße kontrollieren! ---}
\\
\Messbox{100}{50} % Angabe der Breite/Hoehe in mm
\\
\textbf{--- Diese Seite nach dem Druck entfernen! ---}
\end{center}


\end{document}

